\section{その他}
この接で述べることは代数学入門以降の内容である.
\subsection{p群}
$p$群の顕著な性質は以下に代表される.
\begin{prop}
$p$群$G$の中心は非自明である.
\end{prop}
\prf
類等式を考えて, 中心の位数が1だとすると, 中心に属さない元の共役類の位数がp冪で, 類等式の$\bmod{p}$の値で矛盾が起こる.
\qed

$p$群に関する一般的な問題に対しては中心の非自明性に必ず注意されたい. 具体的な$p$群の問題については, そもそも位数の小さい$p$群自体限られているものであるから, いっそのこと$p$群の例を覚えておくのがよい.
\begin{eg} $p$は素数とする.
\ben
\item 位数$p$の群は巡回群$C_p$に限る.
\item 位数$p^2$の群は必ずAbel群であり, 構造定理から$C_{p^2}$あるいは$C_p\times C_p$である. すべての元が位数$p$であると分かれば, $C_p$である.
\item 位数$p^3$のAbel群は$C_{p^3}, C_{p^2}\times C_p, C_p^3$の3通り. 非アーベル群については次の2つがある.
\lefba{($C_{p^2}\rtimes C_p$)
$C_{p^2}^{\times}$には位数$p$の元$a$が存在する. そこで, $\phi: C_p\to \mr{Aut}(C_{p^2})$を$\phi(x) = (\phi_x:t\mapsto a^x t)$により定める. $a^{pm+x} = a^x$だから, これはwell-definedである. これにより$C_{p^2}$の$C_{p}$による半直積$C_{p^2} \rtimes_{\phi} C_{p}$が得られる. その演算は,
\[(n_1,h_1)(n_2,h_2) = (n_1+a^{h_1}n_2, h_1+h_2)\]
というものになる.
}
\lefba{(\tf{Heisenberg群}; $\mr{Heis}(\mb{F}_{p})$)
$\phi:C_p\to \mr{Aut}(C_p\times C_p)$を$\phi(t) = (\phi_t:(a,b)\mapsto (a+bt, b))$により定める($\phi_t$が自己同型であることを確かめよ). これにより半直積$(C_p\times C_p)\rtimes_{\phi} C_p$を得る. この演算は,
\bealn{
&((a_1,b_1),c_1)((a_2,b_2),c_2) \\
&= ((a_1,b_1)+\phi_{c_1}(a_2,b_2), c_1+c_2) \\
&= ((a_1,b_1) + (a_2+c_1b_2, b_2), c_1+c_2) \\
&= ((a_1+a_2+c_1b_2, b_1+b_2), c_1+c_2)
}
そして, この群はHeisenberg群$\mr{Heis}(\mb{F}_{p})$に同型である. この群は, $\mr{GL}_{3}(\mb{F}_{p})$の部分群として, よく次の形で記述される:
\[\mr{Heis}(\mb{F}_{p}) := \set{\bmat{1&c&a\\ 0&1&b\\ 0&0&1} \mid a,b,c\in \mb{F}_{p}}\subset \mr{GL}_{3}(\mb{F}_p)\]
この群では演算が
{\small
\[\bmat{
1&c_1&a_1 \\ 0&1&b_1\\ 0&0&1
}
\bmat{
1&c_2&a_2 \\ 0&1&b_2\\ 0&0&1
}
= \bmat{1&c_{1}+c_2 & a_1 + c_1b_2 + a_2 \\
0&1& b_1+b_2\\ 0&0&1}
\]
}
であり, 先に構成した$(C_p\times C_p)\rtimes_{\phi}C_{p}$と同型であることは視覚的に明らかである. \tf{特に, この群はすべての元が位数$p$である. }
}
\item ただし, 上二つで挙げた群は$p=2$のときに二面体群$D_4$に同型で, 位数$8$の非アーベル群はこのほかに$Q_8 = \set{\pm 1,\pm i, \pm j ,\pm k}$が存在する.
\een
\end{eg}

\begin{prac}
(3)(a)の例において, $a$の取り方には依存するが, $a'$を別の位数$p$の元として取り替えて半直積を構成した場合, それらは同型になることを示せ. (未検証)
\end{prac}


\begin{egprob} (H21数学II-2)
$p$を奇素数とする. $G$が位数$p^3$の有限群で, 単位元以外の各元の位数$p$であるようなものとする. このとき, $G$は$\mr{GL}_{2}(\mb{C})$の部分群と同型ではないことを示せ.
\end{egprob}
\begin{egans}
部分群$H$と同型とせよ. $X\in \mr{GL}_{2}(\mb{C})$で$p$乗して単位行列になるものについて考えよう. 明らかに$X$の固有値は$\zeta_p^{k}$と表せる. もし$X$が対角化可能なら, $X$は対角成分の元が$\zeta_{p}^k$であるものと相似なので, $X^{p} = I$が従う. 逆に$X^p = I$であるなら, $T^p - 1$は重根を持たないので, 最小多項式もそうである. よって$X$が対角化可能になる.
$P\in \mr{GL}_{2}(\mb{C})$のとき,  $H\cong P^{-1}HP$なので$H$は単位行列以外の対角行列を含むとしてよい. それを$A=\mr{diag}(\zeta_p^i, \zeta_p^j)$とする. ここで$\zeta_p^j \neq 1$なら, $kj\equiv 1\pmod{p}$なる$k$を取って$A^k$を考えて取り変えることで$\zeta_p^j = \zeta_p$としてもよい. もし$\zeta_p^j = 1$なら, $\zeta_p^i \neq 1$だから, $S = \sumib{0&1\\ -1&0}$として$H\cong S^{-1}HS$に注意することで$\zeta_p^j \neq 1$の場合に帰着する. 結局, $A=\mr{diag}(\zeta_p^i, \zeta_p)\in H$であると仮定してよい.

$Z(H)$は非自明なので, この中心を解析する. $B=\sumib{a&b\\c&d}\in Z(H)$とせよ.
$AB = \sumib{\zeta_p^i a& \zeta_p^i b\\ \zeta_p c & \zeta_p d}$, $BA = \sumib{\zeta_p^i a & \zeta_p b \\ \zeta_p^i c & \zeta_p d}$だから$i=1$であるか, $i\neq 1$で$b=c=0$である. \par
いずれの場合も, $Z(H)$は対角行列から成る. 単位元でないものが取れるので,そのひとつを$B = \mr{diag}(\zeta_p^{k}, \zeta_p^{l})$とする. \par
位数$p$の対角行列は$p^2$個しかないので, 対角行列ではない$C\in H$が必ず存在する.  $C=\sumib{s&t\\ u&v}\in H$ で$CB = BC$を解くと\ao{$k=l$でなければならないことが分かる.} よって$B_1=\mr{diag}(\zeta_p,\zeta_p) \in Z(H)$となることも分かる. $B$は$Z(H)$の単位元でない任意の元であったから, $Z(H)$は$B_1$の冪からなり, 特に$\sharp Z(H) = p$が分かる. $H$を$P^{-1}HP$に取り替えるにあたって, この中心は変わらないことに注意すると, もともと$B_1$の冪が$H$に含まれていることが分かる. 先の$C$を対角化する行列$Q$を取り$Q^{-1}CQ\in Q^{-1}HQ$を考える. $Q^{-1}CQ$はこのとき$B_1$の冪にはならない対角行列である. よって$H$に$C_0=\mr{diag}(\zeta_p^{k}, \zeta_{p}^{l})\notin Z(H)$が含まれるとしてよい. $B_0$と$C_0$をいくつか掛けていくことで, 任意の位数$p$の対角行列を作ることができる. $K$をこの対角行列全体のなす部分群とする. $C = \sumib{a&b\\ c&d}$が対角行列でないとすれば, $T_{k}=\mr{diag}(\zeta_p^{i_k}, \zeta_p^{j_k})\in K$に対し$T_{1}CT_{2} = \sumib{\zeta_{p}^{i_1+i_2}a & \zeta_{p}^{i_1+j_2}b \\ \zeta_{p}^{i_2+j_1}c & \zeta_{p}^{j_1+j_2}d}$である. $C$は対角行列でないが, $a=d=0$では$C^p \neq I$なので$C$の成分で0のものは1個以下である. このとき, $i_1,i_2,j_1,j_2$を$0$から$p-1$以下で動かすとき, $T_1CT_2$が$p^3$個以上の異なる(対角でない)行列を作ることが分かるが, 単位行列と合わせると$H$の位数を越えるので, 矛盾である. \qed
\end{egans}

\begin{egans} (他の方から教えて貰った有限群の表現論による解法)
埋め込み$\rho:G\to \mr{GL}_{2}(\mb{C})$を2次元表現と見る. 次を認める.
\begin{prop}
一般の群$G$について, 既約$G$表現の次元は$|G|$の約数である.
\end{prop}
これによると$2$は$p^3$の約数にならないので$\rho$は既約表現になり得ない. よって可約であり, $\rho$は自明表現二つの直和に同型で, $\rho$は$\mb{C}^{\times}\times \mb{C}^{\times}$の部分群に同型である. それは$\cyc{p}\times \cyc{p}$の部分群に同型であり, 矛盾. \qed

\end{egans}

\begin{egprob}
(位数$p^3$の群の分類を用いずに)位数$p^3$の群$G$は$|Z(G)| \neq p$を満たすことを示せ.
\end{egprob}
\begin{egans}
$|Z(G)=p^2$であるとしよう. $G/Z(G)$は巡回群であり, これの生成元を$b$とする. このとき任意の$G$の元は$b^kz$ ($z\in Z(G)$)で表せる. このような元を二つ用意し, $b^{k_1}z_1, b^{k_2}z_2$とする.
\[b^{k_1}z_1b^{k_2}z_2 = b^{k_1+k_2}z_1z_2\]
だから$G$はアーベル群である. $|Z(G)| = p^3$なので, 矛盾. よって$|Z(G)|=p^2$ではない.
\end{egans}




\subsection{可解群}
R2の問題では可解群が登場した.
\begin{defn}
$G$が可解群であるとは, 部分群の列$\set{1} = H_0 \subsetneq H_1 \subsetneq \dots \subsetneq H_{n-1} \subsetneq H_n = G$が存在して,
\ben
\item $H_i \norgr H_{i+1}$ ($0\leq i < n-1$)
\item $H_{i+1}/H_i$はAbel群である.
\een
\end{defn}

\begin{rem}
定義から$H_1$は部分Abel群でなければならない.
\end{rem}

正規部分群で繋げなければならないので, 中心$Z(G)$や指数2の部分群であるなど, 一般的に正規部分群と分かる群のことを抑えておくのがよいだろう.


\begin{theo}
$p$群は可解群である.
\end{theo}
\prf $G$の位数$p^n$として$n$に関する帰納法で調べる. $n=1$の場合は明らか. $n-1$までの主張を仮定する.  $Z(G)$は非自明で$G$の正規部分群だから, $G/Z(G)$は$p$群で位数$p^{n-1}$以下. よって$G/Z(G)$は帰納法の仮定より可解. $Z(G)$は可換だから可解.
\qed


\tbox{
\begin{point}
\ben
\item \tf{Sylow $p$ 部分群$P$を考察せよ.} もしそれが正規($n_p = 1$)なら, まず$P\norgr G$. もしくは$G/P$が正規と言えればよい.
\item $N$は正規部分群で, $N$と$G/N$がともに可解群であるようなものが存在するとする. このとき$G$も可解群である.
\een
\end{point}
}{title=テクニック}


\begin{egprob}
位数1998の群$G$は可解群であることを示せ.
\end{egprob}
\begin{egans}
$1998=2 \times 3^3 \times 37$. Sylowの定理よりSylow 37 部分群$S_{37}$が一つあり, それは正規. よって$G/S_{37}$が可解群と言えばよい. $G' := G/S_{37}$のSylow $3$部分群$S_{3}$は, 指数2なので正規部分群. よって$G'/S_3$は明らかに可解群. よって$G$は可解群.
\end{egans}

\begin{prac}
位数$2p^nq^m$ ($p,q$は素数, $2p^n < q$) の群は可解群であることを示せ.
\end{prac}

\begin{egprob}
位数275の群は可解群であることを示せ.
\end{egprob}
\begin{egans}
$275 = 5^2 \times 7 \times 13$. $S_{13}$は可換かつ正規なので$G' = G/S_3$の可解と言えばよい. $G'$のSylow 5部分群$S_{5}$は正規かつ可換(位数$p^2$の群は可換)で $G'/S_{5}$も位数7で可換なので$G'$は正規. よって$G$は正規.
\end{egans}

\begin{egprob}
位数2020の群$G$は可解群であることを示せ.
\end{egprob}
\begin{egans}
$2020 = 4\times 5\times 101$である. $S_{101}$は正規で, 可換より可解なので$G' = G/S_{101}$が可解であればよい. $G'$の$S_{5}$は正規で, 可換より可解なので$G'/S_{5}$が可解であればよい. これは位数$2^2$なので可換であり, 可解である.
\end{egans}

\begin{egprob}
$p,q$が相異なる素数であるとき, 位数$pq$の群$G$は可解であることを示せ.
\end{egprob}
\begin{egans}
$p<q$としてよい. $S_q$は正規で, 可換なので可解である. よって$G/S_{q}$が可解と言えばよいが, これは位数$p$なので可換とくに可解である. よって$G$も可解である. \qed
\end{egans}

次は一筋縄ではいかない場合である. 院試レベルとしては次が出来るようになるのがよい.
\begin{egprob}
$p$が素数であるとき, 位数$4p$の群$G$は可解群であることを示せ.
\end{egprob}
\begin{miss}
$p\geq 5$のとき, Sylow $p$部分群$S_p$は正規でAbel群なので$G':=G/S_p$が可解であればよい. $G'$は位数4だから, これもAbel群でありよい. \par
$p=3$のとき, $S_3$が正規なら先と同様にすればよい. 正規でないこともある(\tf{ここで諦めてはならない！！！}).  正規でないとせよ. このとき$n_3 = 4$であり, 位数3の元が$8$個以上ある. このときSylow 2部分群$S_2$が正規にならねばならない. [証明] もし$n_2 > 1$なら$n_2 = 3$であり, \aka{位数2または4の元が$9$個以上ある}ことになり, $8+9  > 12$だから矛盾. [証明終] よって$S_2$は正規かつ(位数$p^2$だから)可換で可解. そして, $G/S_2$は位数3だから可解. よって$G$も可解. \par
$p=2$のとき, $G$はp群であるから可解である. あるいは, 次のように直接調べる: $Z(G)$が非自明である. $Z(G) = G$なら$G$は可換なのでよい. $Z(G)$の位数2,4なら$Z(G)$は可換で可解かつ$G/Z(G)$の位数が4,2なので可換で可解. よっていずれの場合も可解.
\end{miss}
\begin{egans}
赤線の部分は誤りである(一方で位数3の元が8個あることは正しい！). 9個以上とまではいえない. なぜなら,$S_{2,1},S_{2,2}$がSylow-2部分群, つまり位数4のとき, $S_{2,1}\cap S_{2,2} \neq \set{e}$となるかもしれないからだ. \par
そこで次のように考えるとよい. $S_{2,i}$ ($i=1,2,3$)がSylow-2部分群とする. このとき様々な組み合わせを考えると$S_{2,1}\cup S_{2,2}\cup S_{2,3}$が7個以上の元を持つことが分かり, 位数1,2,4の元が7個以上あることが分かり, 矛盾である. \qed
\end{egans}

\begin{egprob}
位数$2p^n$の群は可解であることを示せ.
\end{egprob}
\begin{egans}
$S_p$は指数2なので正規で可解よりすぐわかる.
\end{egans}

\begin{egprob}
位数$2pq$ ($3\leq p<q$)の群は可解であることを示せ.
\end{egprob}
\begin{egans}
$n_q = 1$の場合は$S_q$が正規で, 位数$2p$の群は可解なのでよい. $n_q\equiv 1\pmod{q}$かつ$n_q \in \set{1,2,p,2p}$であるが, $n_q \neq 1$だとすると$n_q \geq q-1 >  p$なので$n_q = 2p$である. この場合$2p-1 = qk$と書けるが$p<q$により$k=1$である. よって$2p=q+1$となる. \par
$n_q = 2p$のとき位数$q$の元が$2p(q-1)$個だけあり, そうでないものが$2p$個あると言える(この中に単位元もある).
一方$n_p \in \set{1,2,q,2q}$である. $n_p = 1$のときは$G/S_p$が可解なのでよい. $n_p > 1$のときを考える. $n_p(p-1)$個の元が位数$p$であることが分かるから, $n_p(p-1) \leq 2p-1$でなければならない. よって$n_p = 2$である. このとき位数$p$の元は$2p-2$個はあるといえる. よって, 位数$p,q$でない元が単位元を除きあと1個ある. \par
一方$n_2 \in \set{1,p,q,pq}$であり, $n_2 > 1$だと位数2の元が$2$個以上あることになるので現在の状況に矛盾する. よって$n_2=1$で$S_2$が正規. $G/S_2$は可解なのでよい. \qed
\end{egans}
\begin{egprob}
位数$56$の群が可解であることを示せ.
\end{egprob}
\begin{egans}
$n_7$の可能性は1と8. $n_7 = 1$ならよい. $n_7 = 8$なら48個以上, 位数7の元が存在することが保証される. もしこのとき$n_2 = 7$なら位数2の元が$8$個以上(実際はもっと)あることが保証される. これと単位元をあわせると$56$を越えるのでおかしい. よって$n_2 = 1$なので$G/S_2$が可解と言えばよく. これは位数7なのでOK. \qed
\end{egans}




\begin{egprob}
位数$p^2q$ ($p<q$)の群は可解であることを示せ.
\end{egprob}
\begin{egans}
$n_q=1$のときは$G/S_q$が位数$p^2$で可換. $n_q>1$とする. $p-1$は$q$で割れないので, $p^2-1 = (p-1)(p+1)$が$q$で割り切れ, $p+1$が$q$で割り切れなければならない. $p < q \leq p+1$となる. よって$q=p+1$だが$p,q$が素数なので$p=2,q=3$となる. この場合は, $4p$のときの考察で完了している. \qed
\end{egans}

次も知識として知っておくとよい.
\begin{theo}
位数59以下の群は可解群である.
\end{theo}
\prf まず, $1,\dots,59$を次のパターンに分類する.
\ben
\item 素数の累乗: 1, 2, 3, $2^2$, 5, 7, $2^3$, $3^2$, 11, 13, $2^4$, 17, 19, 23, $5^2$, $3^3$ ,29,31, $2^5$ ,37, 41, 43, 47,$7^2$, 53, 59. このときは$p$群なのでよい.
\item 異なる二つの素数の積: $6,10,14,15,21,22,26,33,35,38,39,46$, $51,55,57,59$.
\item $4p$: $12,20,28,36,44,52,$
\item $2p^n$: $18,50,54$
\item $2pq$: $30,42$
\item $p^2q (p<q)$: 45がこれに当てはまるのでよい.
\item それ以外: $24,40,48,56$
\een
24,56についてはすでに調べた.
このうち, 位数が40のものが可解であることは簡単にわかる.
\begin{prac}
位数40の群が可解であることを示せ.
\end{prac}
%8 times 5で n_5 = 1なのでOK
少し難易度が上がるが24も簡単である.
\begin{prac}
位数24の群が可解であることを示せ.
\end{prac}
%\begin{egans} $n_2=1$ならよい. $n_2 = 3$なら位数2冪の元が$21$個はあり, このとき$n_3 \geq 3$にはなりえない. よって$n_3=1$で$G/S_3$が可解なのでよい.
%\end{egans}
よって次を示せばよいが, これは難易度が高い.
\begin{claim}
位数48の群は可解である.
\end{claim}
\prf $S_2$が正規の場合はよい. $S_2$が正規でないとき, $n_2 = 3$であり, それらを$S_{2,i}$ ($i=1,2,3$)とする. ここで\tf{シロー部分群は互いに共役であるから}$N:=S_{2,1}\cap S_{2,2} \cap S_{2,3}$は正規部分群である. もし$N$の位数が2以上なら, $G/N$が位数24,12,6,3のどれかなので$G$も可解であることが分かる. $N=\set{1}$とする. この場合, 実は異なる$i,j$について$S_{2,i}\cap S_{2,j} = \set{1}$であることを示そう. $i=1,j=2$としてよい. $x\in S_{2,1}\cap S_{2,2}$とする. このとき\bolm{x^{-1}S_{2,3}x = S_{2,k}}となる$k$があるが, これは$k=3$でなければならない. よって$x\in N_{G}(S_{2,3})$である. ところで, Sylowの定理より$N_{G}(S_{2,3})$の指数が$n_2 = 3$で, 正規化群の定義より$N_{G}(S_{2,3})\supset S_{2,3}$なので, もし$x\neq 1$なら$x\notin S_{2,3}$だが, このとき$N_{G}(S_{2,3})\supsetneq S_{2,3}$となるのに,  $S_{2,3}$と$N_{G}(S_{2,3})$がともに指数3であるのはおかしい. よって矛盾で, $x=1$だから$S_{2,i}\cap S_{2,j} = \set{1}$である. よって, $G$には位数2,4,8,16の元が少なくとも$45$個ある. もしこのとき$n_3 \geq 2$なら, 位数3の元が少なくとも4個あり, おかしい. よって$n_3 = 1$であり, このとき$G/S_3$が可解だからよい. \qed
\begin{prac}
上の証明をガロア理論から解釈せよ.
\end{prac}



次の定理は知識用である. 答案には使えないであろうが, 覚えておくといいかもしれない.
\begin{theo} (\tf{Burnside}の定理)
$p,q$を素数とする. 位数$p^m q^n$の群は可解である.
\end{theo}
よって, 位数$2021 = 43\times 47$の群は可解群である. また, 位数$4p$の群はやはり可解である.
\begin{prac}
位数2021の群は可解であることをBurnsideの定理を使わずに示せ.
\end{prac}

最後に, 位数60以降の話を少しだけ見てみよう.
\begin{prop}
位数$\dfrac{n!}{2}$ ($n\geq 5$)の群は可解とは限らない.
\end{prop}
\prf $A_n$($n\geq 5$)は単純群だから. \qed

\begin{egprob}
$S_{n}$ ($n\geq 5$)が可解群ではないことを示せ. $A_n$が単純群であることは用いてよい.
\end{egprob}
\begin{egans} $S_n$の正規部分群$N\neq \set{1}$を考える. このとき$N\cap A_n\norgr A_n$より$N\cap A_5 = \set{1}$でなければならない. どのような部分群も偶置換は必ず含むので, $N=A_n$となるしかない. よって系列は$A_n\norgr S_n$で終わらなければならず, そのあと$A_n$の部分群の列を作らなければならないが, これは$A_n$の単純性から不可能. \qed
\end{egans}


\begin{prac}
位数61,62,63,64,65,66,67,68,69,70の群は可解であることを示せ.
\end{prac}





\subsection{演習問題}





\clearpage
